\input{../tex-inputs/latex-korrekturansicht-vorspann}

\section[Arthur Schnitzler an Gustav Schwarzkopf, 18. 8. 1928]{L04349 Arthur Schnitzler an Gustav Schwarzkopf, 18. 8. 1928}
\nopagebreak\mylabel{L04349v}
\rehead{ }\normalsize\beginnumbering\briefempfaengerindex{Schwarzkopf, Gustav@\textsc{Schwarzkopf, Gustav}!zzzSchnitzler, Arthur@\emph{von Arthur Schnitzler}!1928-08-181@{18. 8. 1928}|(be}
\toendnotes[C]{\smallbreak\pagebreak[2]}
\correspDesc{Versand  durch Arthur Schnitzler am 18. 8. 1928 in Hohenschwangau
\newline{}Erhalt  durch Gustav Schwarzkopf im Zeitraum [19. 8. 1928 – 23. 8. 1928?] in Wien}\toendnotes[C]{\smallbreak}
\Standort{DLA, A:Schnitzler, HS.1985.1.1897.}
\physDesc{Bildpostkarte, 462 Zeichen
\newline{}Handschrift: Bleistift, lateinische Kurrent
\newline{}Versand: Stempel: »\nobreak{}\oindex{Hohenschwangau@\textbf{Hohenschwangau}, \emph{Teil eines besiedelten Ortes}|pwk}Hohens{[}chwangau{]}, \textcolor{gray}{18} Aug 2{[}8{]}\nobreak{}«.  }\toendnotes[C]{\smallbreak}\pstart{}{\pb}Herr Gustav Schwarzkopf\pend{}\pstart{}\textcolor{pink}{Wien I}\oindex{I., Innere Stadt@\textbf{I., Innere Stadt}, \emph{Verwaltungsgebiet}|pw}{}\ledrightnote{\textcolor{pink}{I., Innere Stadt}}\pend{}\pstart{}\textcolor{pink}{Tiefer Graben 17}\oindex{Wien@\textbf{Wien}!I., Innere Stadt@\textbf{I., Innere Stadt}!Tiefer Graben 17@\textbf{Tiefer Graben 17}, \emph{Wohngebäude}|pw}{}\ledrightnote{\textcolor{pink}{Tiefer Graben 17}}\pend{}{\bigskip}
\pstart
           \noindent{}\centering{}{\pb}\textcolor{gray}{\textbf{\textcolor{pink}{Schloss Neuschwanstein}\oindex{Schloss Neuschwanstein@\textbf{Schloss Neuschwanstein}, \emph{Burg}|pw}{}\ledrightnote{\textcolor{pink}{Schloss Neuschwanstein}} – \textcolor{pink}{Hohenschwangau}\oindex{Hohenschwangau@\textbf{Hohenschwangau}, \emph{Teil eines besiedelten Ortes}|pw}{}\ledrightnote{\textcolor{pink}{Hohenschwangau}} u. \textcolor{pink}{Alpsee}\oindex{Alpsee@\textbf{Alpsee}, \emph{See}|pw}{}\ledrightnote{\textcolor{pink}{Alpsee}}.}}\pend
           \vspace{1em}
\pstart
           \raggedleft{}{\pb}18/8 928\pend
           
\pstart
           \raggedleft{}\textcolor{pink}{Hohenschwangau, Hotel 
                     Alpenrose}\oindex{Hotel Alpenrose [Hohenschwangau]@\textbf{Hotel Alpenrose [Hohenschwangau]}, \emph{Hotel}|pw}{}\ledrightnote{\textcolor{pink}{Hotel Alpenrose [Hohenschwangau]}}\pend
           \vspace{0.5em}
\pstart
           An dieser guten Sitte soll sich  nichts ändern – auch von hier aus will ich Ihnen
               viele tägliche Grüße  senden. –  Im Jahre \label{K_L04349-1v}\edtext{1909 war ich hier ein paar Wochen \uline{vorher}}{\lemma{\textnormal{\emph{1909 … vorher}}}\Cendnote{\textnormal{\textcolor{blue}{Schnitzler} bezieht sich auf den Tod seiner
                  Tochter \textcolor{blue}{Lili}\pwindex{Cappellini, Lili 13.\,9.\,1909 Wien – 26.\,7.\,1928 Venedig@\textsc{Cappellini, Lili} (13.\,9.\,1909 Wien – 26.\,7.\,1928 Venedig)|pwk}, die am 13. 9. 1909 – und
                  damit kurz nach seinem ersten Besuch in \textcolor{pink}{Hohenschwangau}\oindex{Hohenschwangau@\textbf{Hohenschwangau}, \emph{Teil eines besiedelten Ortes}|pwk} am 30. 8. 1909 – auf die Welt kam und die am 26. 7. 1928 gestorben
                  war.}}}\label{K_L04349-1} – wie jetzt ein paar Wochen nachher{\dots}
               Dazwischen liegt nicht mehr und nicht weniger als ein Leben. Und alles Glück, was man
                  \introOben{}\textcolor{gray}{davon}\introOben{} gehabt – man bezahlt es in solchen Tagen tausendfach. Auf Wiedersehen\pend
           \pstart I{\geminationm}er Ihr getreuer \spacefill\mbox{AS}\pend{}\selectlanguage{ngerman}\endnumbering\briefempfaengerindex{Schwarzkopf, Gustav@\textsc{Schwarzkopf, Gustav}!zzzSchnitzler, Arthur@\emph{von Arthur Schnitzler}!1928-08-181@{18. 8. 1928}|)be}\mylabel{L04349h}
\begin{anhang}
\end{anhang}\newcommand{\dateiname}{L04349}\newcommand{\titel}{Arthur Schnitzler an Gustav Schwarzkopf, 18. 8. 1928}\newcommand{\editorInnen}{Herausgegeben von Jahnke, SelmaMüller, Martin Anton}\input{../tex-inputs/latex-korrekturansicht-abspann}
